% appendix-b-hyperparameters.tex -- part of thesis_v3.tex; do not add \documentclass here.
% Draft with the thesis-chapter workflow. Unsupported claims must use \fillin{...}.

\chapter{Hyperparameter Configuration for Training and Optimisation}
\label{ch:appendix_hyperparameters}

\section{Scope and how to read these tables}
\label{sec:appb_scope}

This appendix records the configuration files that drive training and fine-tuning of the
surrogate \texttt{exp059\_capacity\_freq}, the fine-tune override files that sit on top of
the base configuration, and the \texttt{CONFIG} block of the Stage-2 latent-optimisation
script. It is a transcription, not an analysis: every key name and every value is
reproduced exactly as it appears in the source file, with no rounding, no unit
conversion and no reformatting.

\textbf{These are the current contents of the configuration files, not a run record.}
\texttt{config.yaml} has been edited since the capacity-phase training run --- at least in
the two keys \texttt{layout\_train\_prob} and \texttt{occ\_only\_encode\_prob}, discussed
as C5 in \autoref{sec:appb_conflicts} --- and the fine-tune configuration that the
active-learning loop actually launches is regenerated at every launch and overwritten
(C1). Nothing in this appendix may be read as a per-run snapshot of the values that were
in force when a particular checkpoint was produced.

The training configuration lives in \texttt{experiments/exp059\_capacity\_freq/config.yaml}.
Despite the extension, the file is a JSON document with 183 top-level keys. It is read by
\texttt{load\_yaml\_config()} (\texttt{experiments/exp059\_capacity\_freq/codes/exp059\_common.py},
lines 17--25), which discards blank lines and every line whose first non-space character
is \texttt{\#}, then calls \texttt{json.loads} on the remainder. Consequently YAML syntax
--- bare keys, \texttt{-} lists, anchors --- is not accepted. The base \texttt{config.yaml}
contains no comment lines at all; the two override files do.

The 183 keys are grouped below into 21 tables. \textbf{The grouping is editorial}: keys
belonging to one table are in general not contiguous in the file, and the file itself
imposes no ordering or sectioning. The grouping was introduced here so that the tables
are readable; it should not be taken as structure present in the artifact.

None of the 183 keys carries a comment or any other in-file documentation, so this
appendix cannot state why an individual value was chosen.
\fillin{for the values that were actually tuned rather than inherited from exp057/exp058,
add the provenance --- the sweep, ablation or hand-tuning step that fixed them.
\texttt{experiments/exp059\_capacity\_freq/notes.md} would be the artifact to check; no
sweep record was located during evidence extraction.}

\section{Surrogate training configuration: \texttt{exp059\_capacity\_freq}}
\label{sec:appb_exp059_config}

\subsection{Latent structure and architecture keys}
\label{subsec:appb_arch}

Five of the architecture-level switches are set explicitly in \texttt{config.yaml}. The
remaining architecture switches are not present in the file and fall back to code
defaults; they are listed separately in \autoref{sec:appb_code_defaults} so that this
table does not imply they were configured.

\begin{table}[htbp]
\centering
\caption{Latent and architecture keys set in \texttt{config.yaml} (5 keys).}
\label{tab:appb_arch}
\begin{tabular}{@{}ll@{}}
\hline
Key & Value \\
\hline
\texttt{latent\_dim} & 128 \\
\texttt{heatmap\_private\_dim} & 40 \\
\texttt{cond\_dim} & 32 \\
\texttt{freq\_fourier\_features} & 16 \\
\texttt{use\_heatmap\_film} & true \\
\hline
\end{tabular}
\end{table}

\subsection{Training schedule, optimiser and learning-rate scheduler}
\label{subsec:appb_schedule}

\autoref{tab:appb_schedule} gives the schedule keys. The optimiser itself is not
configurable from the file: \texttt{train\_core.py} (lines 906--920) constructs
\texttt{torch.optim.AdamW} with \texttt{weight\_decay} $=$ \texttt{1e-4}, first
attempting the fused implementation on CUDA, then the \texttt{foreach} implementation,
and falling back to plain \texttt{AdamW}. The learning-rate schedule is
\texttt{ReduceLROnPlateau(mode="min", factor=c.lr\_factor, patience=c.lr\_patience,
min\_lr=c.lr\_min)} (\texttt{train\_core.py}, lines 1431--1433), which with the values in
\autoref{tab:appb_schedule} resolves to factor $0.5$, patience $15$ and minimum learning
rate \texttt{1e-05}.

On resume with \texttt{reset\_lr\_on\_resume} set, every parameter group's learning rate
is reset to \texttt{learning\_rate}, except parameter group index 1 when the resume epoch
is at or past \texttt{impedance\_peak\_focus\_epoch}, which is set to
\texttt{learning\_rate} $\times$ \texttt{impedance\_decoder\_lr\_mult}; the scheduler is
then recreated from scratch (\texttt{train\_core.py}, lines 1442--1451).

\begin{table}[htbp]
\centering
\caption{Training schedule and scheduler keys (13 keys).}
\label{tab:appb_schedule}
\begin{tabular}{@{}ll@{}}
\hline
Key & Value \\
\hline
\texttt{num\_epochs} & 510 \\
\texttt{batch\_size} & 224 \\
\texttt{val\_batch\_size} & 320 \\
\texttt{learning\_rate} & 1.5e-05 \\
\texttt{lr\_min} & 1e-05 \\
\texttt{lr\_patience} & 15 \\
\texttt{lr\_factor} & 0.5 \\
\texttt{train\_split} & 0.9 \\
\texttt{num\_workers} & 4 \\
\texttt{reset\_lr\_on\_resume} & true \\
\texttt{train\_samples\_per\_epoch} & 64000 \\
\texttt{impedance\_decoder\_lr\_mult} & 1.0 \\
\texttt{impedance\_separate\_lr} & true \\
\hline
\end{tabular}
\end{table}

\subsection{Data splitting and class balancing}
\label{subsec:appb_split}

\begin{table}[htbp]
\centering
\caption{Data split and balancing keys (7 keys).}
\label{tab:appb_split}
\begin{tabular}{@{}ll@{}}
\hline
Key & Value \\
\hline
\texttt{balance\_k} & false \\
\texttt{k\_balance\_power} & 0.4 \\
\texttt{k\_balance\_smoothing} & 0.001 \\
\texttt{stratify\_by\_k} & true \\
\texttt{split\_by\_design} & true \\
\texttt{balance\_freq} & false \\
\texttt{freq\_balance\_power} & 1.4 \\
\hline
\end{tabular}
\end{table}

Note that \texttt{balance\_k} and \texttt{balance\_freq} are both \texttt{false} while
\texttt{k\_balance\_power}, \texttt{k\_balance\_smoothing} and
\texttt{freq\_balance\_power} carry non-default-looking values, i.e.\ the two weighting
exponents are present but inactive under this configuration.
\fillin{confirm from \texttt{train\_core.py} whether \texttt{k\_balance\_power} and
\texttt{freq\_balance\_power} are read at all when the corresponding \texttt{balance\_*}
flag is false, and state the sampler behaviour that is actually used.}

\subsection{Dataset, overlay and checkpoint paths}
\label{subsec:appb_paths}

\begin{table}[htbp]
\centering
\caption{Dataset, overlay and checkpointing keys (8 keys).}
\label{tab:appb_paths}
\begin{tabular}{@{}lp{0.44\textwidth}@{}}
\hline
Key & Value \\
\hline
\texttt{data\_dir} & \texttt{datasets/\allowbreak data\_multifreq\_\allowbreak train\_norm\_unbounded} \\
\texttt{experiment\_dir} & \texttt{experiments/\allowbreak exp059\_capacity\_freq} \\
\texttt{al\_overlay\_data\_dir} & \texttt{datasets/\allowbreak data\_multifreq\_\allowbreak al\_overlay\_exp059} \\
\texttt{al\_overlay\_sample\_weight} & 128.0 \\
\texttt{resume\_checkpoint} & \texttt{experiments/\allowbreak exp059\_capacity\_freq/\allowbreak checkpoints/\allowbreak last\_model.pt} \\
\texttt{recalculate\_curriculum\_on\_resume} & false \\
\texttt{checkpoint\_interval} & 5 \\
\texttt{keep\_last\_n\_checkpoints} & 0 \\
\hline
\end{tabular}
\end{table}

\subsection{Logging and validation cadence}
\label{subsec:appb_logging}

\begin{table}[htbp]
\centering
\caption{Logging and validation cadence keys (3 keys).}
\label{tab:appb_logging}
\begin{tabular}{@{}ll@{}}
\hline
Key & Value \\
\hline
\texttt{val\_on\_checkpoint\_only} & true \\
\texttt{epoch\_print\_interval} & 1 \\
\texttt{epoch\_log\_interval} & 2 \\
\hline
\end{tabular}
\end{table}

\clearpage

\subsection{Heatmap loss weights}
\label{subsec:appb_heatmap_loss}

The heatmap branch carries the largest number of individually weighted loss terms.
\autoref{tab:appb_heatmap_loss} lists the general reconstruction and structure weights,
\autoref{tab:appb_heatmap_region} the peak- and valley-region terms, and
\autoref{tab:appb_heatmap_focus} the keys governing the heatmap focus phase.
\fillin{describe, in Chapter 3 rather than here, what each heatmap loss term penalises
and why the relative weights were set this way; the config file records no rationale.}

\begin{table}[htbp]
\centering
\caption{Heatmap reconstruction and structure loss weights (15 keys).}
\label{tab:appb_heatmap_loss}
\begin{tabular}{@{}ll@{}}
\hline
Key & Value \\
\hline
\texttt{heatmap\_weight} & 5.0 \\
\texttt{heatmap\_phys\_p99\_weight} & 2.0 \\
\texttt{heatmap\_phys\_p99\_percentile} & 99.0 \\
\texttt{heatmap\_peak\_weight} & 3.0 \\
\texttt{heatmap\_grad\_weight} & 1.5 \\
\texttt{heatmap\_lap\_weight} & 2.0 \\
\texttt{heatmap\_contrast\_weight} & 1.5 \\
\texttt{heatmap\_contrast\_margin} & 0.5 \\
\texttt{heatmap\_bg\_weight} & 0.5 \\
\texttt{heatmap\_dynrange\_weight} & 2.0 \\
\texttt{heatmap\_pearson\_weight} & 0.75 \\
\texttt{heatmap\_grad\_vector\_weight} & 2.0 \\
\texttt{heatmap\_grad\_direction\_weight} & 1.0 \\
\texttt{heatmap\_grad\_direction\_min\_mag} & 0.08 \\
\texttt{heatmap\_grad\_huber\_delta} & 0.5 \\
\hline
\end{tabular}
\end{table}

\begin{table}[htbp]
\centering
\caption{Heatmap peak- and valley-region loss terms (11 keys).}
\label{tab:appb_heatmap_region}
\begin{tabular}{@{}ll@{}}
\hline
Key & Value \\
\hline
\texttt{heatmap\_peak\_centroid\_top\_q} & 0.95 \\
\texttt{heatmap\_valley\_centroid\_bottom\_q} & 0.05 \\
\texttt{heatmap\_peak\_topregion\_weight} & 2.5 \\
\texttt{heatmap\_valley\_topregion\_weight} & 2.5 \\
\texttt{heatmap\_peak\_topregion\_q} & 0.95 \\
\texttt{heatmap\_valley\_topregion\_q} & 0.05 \\
\texttt{heatmap\_topregion\_huber\_delta} & 0.5 \\
\texttt{heatmap\_peak\_loc\_weight} & 2.5 \\
\texttt{heatmap\_peak\_loc\_temperature} & 0.035 \\
\texttt{heatmap\_valley\_loc\_weight} & 0.0 \\
\texttt{heatmap\_valley\_loc\_temperature} & 0.03 \\
\hline
\end{tabular}
\end{table}

\texttt{heatmap\_valley\_loc\_weight} is $0.0$, so the valley-localisation term is
switched off while its temperature is still specified.

\begin{table}[htbp]
\centering
\caption{Heatmap focus phase and low-$K$ emphasis (9 keys).}
\label{tab:appb_heatmap_focus}
\begin{tabular}{@{}ll@{}}
\hline
Key & Value \\
\hline
\texttt{heatmap\_focus\_start\_frac} & 0.0 \\
\texttt{heatmap\_focus\_start\_epoch} & 25 \\
\texttt{heatmap\_focus\_heatmap\_weight} & 7.0 \\
\texttt{heatmap\_focus\_impedance\_weight} & 4.0 \\
\texttt{heatmap\_focus\_cross\_freq\_weight} & 1.1 \\
\texttt{heatmap\_focus\_dynrange\_weight} & 3.0 \\
\texttt{heatmap\_focus\_modality\_dropout} & 0.03 \\
\texttt{hm\_low\_k\_threshold} & 3 \\
\texttt{hm\_low\_k\_multiplier} & 2.0 \\
\hline
\end{tabular}
\end{table}

\begin{table}[htbp]
\centering
\caption{Heatmap value range and clipping (4 keys).}
\label{tab:appb_heatmap_clip}
\begin{tabular}{@{}ll@{}}
\hline
Key & Value \\
\hline
\texttt{background\_value} & -3.1792 \\
\texttt{heatmap\_z\_clip\_min} & -1.6791870594024658 \\
\texttt{heatmap\_z\_clip\_max} & 6.275787830352783 \\
\texttt{heatmap\_clip\_recon} & false \\
\hline
\end{tabular}
\end{table}

\texttt{heatmap\_z\_clip\_min} and \texttt{physics\_fg\_clip\_min}
(\autoref{tab:appb_physics}) hold the same value, $-1.6791870594024658$.
\fillin{state where this number comes from --- presumably a percentile of the normalised
heatmap distribution recorded in \texttt{normalization\_stats.json} --- and confirm that
the two keys are meant to be tied. Do not keep the word "presumably": either cite the
statistics file or mark the link as unverified.}

\clearpage

\subsection{Impedance loss weights and peak curriculum}
\label{subsec:appb_impedance_loss}

The impedance branch predicts the spectrum $Z(f)$ over the 231 frequency bins spanning
1--600\,MHz. \autoref{tab:appb_impedance} lists its loss weights;
\autoref{tab:appb_imp_curriculum} lists the curriculum keys that stage the peak-alignment
terms during training.

\begin{table}[htbp]
\centering
\caption{Impedance loss weights (11 keys).}
\label{tab:appb_impedance}
\begin{tabular}{@{}ll@{}}
\hline
Key & Value \\
\hline
\texttt{impedance\_weight} & 4.0 \\
\texttt{impedance\_deriv\_weight} & 1.4 \\
\texttt{impedance\_topk\_k} & 20 \\
\texttt{impedance\_topk\_weight} & 4.5 \\
\texttt{impedance\_under\_penalty} & 2.8 \\
\texttt{impedance\_concavity\_weight} & 2.5 \\
\texttt{impedance\_freq\_weight\_alpha} & 2.0 \\
\texttt{impedance\_dual\_topk\_weight} & 0.75 \\
\texttt{impedance\_peak\_index\_weight} & 1.5 \\
\texttt{impedance\_peak\_mag\_weight} & 2.0 \\
\texttt{impedance\_num\_peaks} & 8 \\
\hline
\end{tabular}
\end{table}

\begin{table}[htbp]
\centering
\caption{Impedance peak curriculum (6 keys).}
\label{tab:appb_imp_curriculum}
\begin{tabular}{@{}ll@{}}
\hline
Key & Value \\
\hline
\texttt{impedance\_peak\_start\_frac} & 0.1 \\
\texttt{impedance\_peak\_ramp\_frac} & 0.05 \\
\texttt{impedance\_peak\_focus\_frac} & 0.1 \\
\texttt{impedance\_peak\_start\_epoch} & 25 \\
\texttt{impedance\_peak\_ramp\_epochs} & 22 \\
\texttt{impedance\_peak\_focus\_epoch} & 25 \\
\hline
\end{tabular}
\end{table}

Each curriculum stage is specified twice, once as a fraction of the run and once as an
absolute epoch.
\fillin{state which of the two forms \texttt{train\_core.py} actually uses --- the
\texttt{*\_frac} keys or the \texttt{*\_epoch} keys --- and whether
\texttt{recalculate\_curriculum\_on\_resume} (false, \autoref{tab:appb_paths}) selects
between them on resume.}

\subsection{Occupancy loss}
\label{subsec:appb_occupancy_loss}

The occupancy head predicts the placement vector $\mathbf{b} \in \{0,1\}^{52}$ under a
budget $K$.

\begin{table}[htbp]
\centering
\caption{Occupancy loss keys (11 keys).}
\label{tab:appb_occupancy}
\begin{tabular}{@{}ll@{}}
\hline
Key & Value \\
\hline
\texttt{occupancy\_weight} & 7.0 \\
\texttt{occupancy\_focal\_gamma} & 1.5 \\
\texttt{focal\_gamma\_warmup\_frac} & 0.0 \\
\texttt{focal\_gamma\_warmup\_epochs} & 0 \\
\texttt{occ\_k\_consistency\_weight} & 2.0 \\
\texttt{use\_k\_weighting} & true \\
\texttt{occ\_mid\_k\_center} & 25.0 \\
\texttt{occ\_mid\_k\_sigma} & 8.0 \\
\texttt{occ\_mid\_k\_boost} & 0.75 \\
\texttt{occupancy\_binary\_decode} & true \\
\texttt{occupancy\_binary\_ste\_train} & false \\
\hline
\end{tabular}
\end{table}

\subsection{KL and latent regularisation}
\label{subsec:appb_kl}

\begin{table}[htbp]
\centering
\caption{KL, $\beta$-annealing and latent regularisation keys (17 keys).}
\label{tab:appb_kl}
\begin{tabular}{@{}ll@{}}
\hline
Key & Value \\
\hline
\texttt{free\_bits} & 0.12 \\
\texttt{mu\_hinge\_threshold} & 4.0 \\
\texttt{mu\_hinge\_weight} & 0.1 \\
\texttt{mu\_bias\_weight} & 0.05 \\
\texttt{per\_expert\_kl\_weight} & 0.02 \\
\texttt{sigma\_reg\_weight} & 1.5 \\
\texttt{sigma\_reg\_target} & 0.45 \\
\texttt{use\_beta\_annealing} & false \\
\texttt{beta\_start\_epoch} & 0 \\
\texttt{beta\_end\_frac} & 0.3 \\
\texttt{beta\_initial} & 0.0 \\
\texttt{beta\_final} & 0.03 \\
\texttt{beta\_phase2\_final} & 0.03 \\
\texttt{beta\_phase2\_end\_frac} & 0.3 \\
\texttt{beta\_end\_epoch} & 25 \\
\texttt{beta\_phase2\_end\_epoch} & 25 \\
\texttt{kan\_spline\_l1\_weight} & 0.0001 \\
\hline
\end{tabular}
\end{table}

\texttt{use\_beta\_annealing} is \texttt{false}, so the eight $\beta$ schedule keys
describe a schedule that is not active; \texttt{beta\_final} is $0.03$.
\fillin{confirm from \texttt{train\_core.py} which single $\beta$ value is applied when
\texttt{use\_beta\_annealing} is false.}

\clearpage

\subsection{Modality dropout and cross-modal terms}
\label{subsec:appb_moddrop}

Modality dropout is what makes the occupancy-only encode path usable at inference, which
is the only path available when a placement is proposed rather than reconstructed.

\begin{table}[htbp]
\centering
\caption{Modality dropout and cross-modal keys (8 keys).}
\label{tab:appb_moddrop}
\begin{tabular}{@{}ll@{}}
\hline
Key & Value \\
\hline
\texttt{modality\_dropout} & 0.06 \\
\texttt{modality\_dropout\_start} & 0.06 \\
\texttt{modality\_dropout\_anneal\_frac} & 0.04 \\
\texttt{modality\_dropout\_anneal\_epochs} & 18 \\
\texttt{modality\_dropout\_protect\_heatmap\_frac} & 0.325 \\
\texttt{modality\_dropout\_protect\_heatmap\_epoch} & 25 \\
\texttt{cross\_modal\_weight} & 0.75 \\
\texttt{cross\_modal\_update\_freq} & 4 \\
\hline
\end{tabular}
\end{table}

\subsection{Cross-frequency and layout path}
\label{subsec:appb_crossfreq}

\begin{table}[htbp]
\centering
\caption{Cross-frequency and layout-path keys (11 keys).}
\label{tab:appb_crossfreq}
\begin{tabular}{@{}ll@{}}
\hline
Key & Value \\
\hline
\texttt{cross\_freq\_weight} & 1.2 \\
\texttt{cross\_freq\_start\_frac} & 0.025 \\
\texttt{cross\_freq\_start\_epoch} & 20 \\
\texttt{cross\_freq\_layout\_z\_only} & false \\
\texttt{cross\_freq\_layout\_mix\_prob} & 0.45 \\
\texttt{cross\_freq\_use\_encode\_z} & true \\
\texttt{layout\_train\_prob} & 0.65 \\
\texttt{occ\_only\_encode\_prob} & 0.4 \\
\texttt{eval\_use\_occ\_only\_layout} & true \\
\texttt{freq\_jitter\_prob} & 0.6 \\
\texttt{freq\_jitter\_log10\_sigma} & 0.1 \\
\hline
\end{tabular}
\end{table}

The values of \texttt{layout\_train\_prob} and \texttt{occ\_only\_encode\_prob} in this
table are the current contents of the file, not a record of the capacity-phase run.
\texttt{docs/cursor-handoff.md} states that the capacity phase was run with
\texttt{layout\_train\_prob} $=0$ and \texttt{occ\_only\_encode\_prob} $=0$, whereas
\texttt{config.yaml} now holds $0.65$ and $0.4$. The file therefore no longer describes
the run that produced the capacity checkpoint, and a reader must not treat
\autoref{tab:appb_crossfreq} as the capacity-phase record for these two keys.
\fillin{recover the two values that were in force when the capacity checkpoint was
trained --- the training log or the config snapshot stored alongside
\texttt{checkpoints/last\_model.pt} would settle it --- and state whether any other key
in \autoref{sec:appb_exp059_config} was also edited after that run.}

\subsection{Distillation and physics terms}
\label{subsec:appb_distill_physics}

\begin{table}[htbp]
\centering
\caption{Latent- and output-distillation keys (4 keys).}
\label{tab:appb_distill}
\begin{tabular}{@{}ll@{}}
\hline
Key & Value \\
\hline
\texttt{latent\_distill\_weight} & 2.5 \\
\texttt{latent\_distill\_private\_mult} & 3.0 \\
\texttt{latent\_distill\_logvar\_weight} & 0.3 \\
\texttt{output\_distill\_weight} & 2.5 \\
\hline
\end{tabular}
\end{table}

\begin{table}[htbp]
\centering
\caption{Physics-loss keys (10 keys).}
\label{tab:appb_physics}
\begin{tabular}{@{}ll@{}}
\hline
Key & Value \\
\hline
\texttt{physics\_ri\_weight} & 1.0 \\
\texttt{physics\_critic\_sup\_weight} & 2.0 \\
\texttt{physics\_ar\_weight} & 0.5 \\
\texttt{physics\_fg\_clip\_min} & -1.6791870594024658 \\
\texttt{physics\_critic\_warmup\_frac} & 0.025 \\
\texttt{physics\_slope\_anneal\_frac} & 0.05 \\
\texttt{physics\_critic\_warmup\_epochs} & 11 \\
\texttt{physics\_slope\_anneal\_epochs} & 22 \\
\texttt{penalty\_warmup\_frac} & 0.025 \\
\texttt{penalty\_warmup\_epochs} & 11 \\
\hline
\end{tabular}
\end{table}

\clearpage

\subsection{Off-anchor evaluation and fine-tune early stopping}
\label{subsec:appb_eval_keys}

The \texttt{eval\_off\_anchor} keys define the frequencies at which the surrogate is scored
outside the normal validation pass. The name is misleading and the set must not be read as
a held-out frequency set. Of the three frequencies in \texttt{eval\_off\_anchor\_mhz},
$155.0$ and $265.0$\,MHz are not training anchors, but $250.0$\,MHz is one of the 24
anchors listed in \texttt{datasets/data\_multifreq\_train\_norm\_unbounded/\allowbreak
dataset\_meta.json} (\texttt{pi\_frequencies\_mhz}) and reproduced in
\autoref{ch:appendix_simulation_dataset}. It also carries the largest weight of the three,
$3.0$ against $2.5$ for each of the other two. One third of this evaluation set is
therefore trained on, and it is the most heavily weighted third.

This matters beyond bookkeeping, because the metric is also a checkpoint-selection
criterion: \texttt{experiments/exp059\_capacity\_freq/checkpoints/\allowbreak
best\_off\_anchor\_model.pt} exists, so a checkpoint is retained on the basis of a score
computed partly at a trained frequency. Any statement in Chapter 4 or Chapter 5 about
generalisation to unseen frequencies must be made against $155.0$ and $265.0$\,MHz alone,
not against the pooled off-anchor number.
\fillin{state why 250.0\,MHz, an existing training anchor, is included in the off-anchor
evaluation set, and whether the off-anchor metrics reported elsewhere in the thesis are
per frequency or pooled across the three; the off-anchor evaluation code and the metrics
CSVs under \texttt{experiments/exp059\_capacity\_freq/metrics/} would settle it.}

\begin{table}[htbp]
\centering
\caption{Off-anchor evaluation keys (4 keys). \texttt{250.0} is a training anchor; see the
text above.}
\label{tab:appb_offanchor}
\begin{tabular}{@{}lp{0.42\textwidth}@{}}
\hline
Key & Value \\
\hline
\texttt{eval\_off\_anchor\_mhz} & \texttt{[ 155.0, 250.0, 265.0 ]} \\
\texttt{eval\_off\_anchor\_mhz\_weights} & \texttt{\{ "155": 2.5, "250": 3.0, "265": 2.5 \}} \\
\texttt{eval\_off\_anchor\_max\_batches} & 20 \\
\texttt{eval\_off\_anchor\_interval} & 5 \\
\hline
\end{tabular}
\end{table}

\begin{table}[htbp]
\centering
\caption{Active-learning fine-tune early-stopping keys present in the base
\texttt{config.yaml} (4 keys).}
\label{tab:appb_earlystop}
\begin{tabular}{@{}ll@{}}
\hline
Key & Value \\
\hline
\texttt{al\_finetune\_early\_stop} & true \\
\texttt{al\_finetune\_early\_stop\_patience} & 3 \\
\texttt{al\_finetune\_early\_stop\_min\_delta} & 0.01 \\
\texttt{al\_finetune\_early\_stop\_kind} & \texttt{"layout\_cross"} \\
\hline
\end{tabular}
\end{table}

\subsection{Numerical stability}
\label{subsec:appb_numerics}

\begin{table}[htbp]
\centering
\caption{Numerical-stability keys (10 keys).}
\label{tab:appb_numerics}
\begin{tabular}{@{}ll@{}}
\hline
Key & Value \\
\hline
\texttt{training\_force\_fp32} & true \\
\texttt{training\_sanitize\_grads} & false \\
\texttt{training\_skip\_nonfinite\_batches} & true \\
\texttt{training\_grad\_clip\_norm} & 1.0 \\
\texttt{training\_logvar\_clamp\_min} & -10.0 \\
\texttt{training\_logvar\_clamp\_max} & 10.0 \\
\texttt{training\_recon\_z\_soft\_floor} & -8.0 \\
\texttt{training\_recon\_z\_soft\_cap} & null \\
\texttt{training\_recon\_z\_cap\_mult} & 1.05 \\
\texttt{training\_loss\_finite\_cap} & 1000000.0 \\
\hline
\end{tabular}
\end{table}

\subsection{Hardware and runtime}
\label{subsec:appb_runtime}

\begin{table}[htbp]
\centering
\caption{Hardware and runtime keys (12 keys).}
\label{tab:appb_hardware}
\begin{tabular}{@{}ll@{}}
\hline
Key & Value \\
\hline
\texttt{amp} & \texttt{"off"} \\
\texttt{tf32} & false \\
\texttt{compile} & false \\
\texttt{compile\_mode} & \texttt{"reduce-overhead"} \\
\texttt{cache\_in\_ram} & false \\
\texttt{use\_ddp} & true \\
\texttt{ddp\_base\_batch\_size} & 160 \\
\texttt{ddp\_linear\_lr\_scale} & true \\
\texttt{persistent\_workers} & false \\
\texttt{prefetch\_factor} & 2 \\
\texttt{empty\_cache\_interval} & 0 \\
\texttt{device} & \texttt{"cuda:0"} \\
\hline
\end{tabular}
\end{table}

Mixed precision is off and \texttt{training\_force\_fp32} is true, so training ran in
single precision. \texttt{use\_ddp} is true with \texttt{ddp\_base\_batch\_size} $=160$,
\texttt{batch\_size} $=224$ and \texttt{ddp\_linear\_lr\_scale} $=$ true.
\fillin{state the effective per-GPU and global batch size and the learning rate that
linear scaling actually produced. No artifact recording the scaled learning rate was
found; a training log line from the DDP run would settle it.}

\clearpage

\section{Architecture switches taken from code defaults}
\label{sec:appb_code_defaults}

Not every architecture switch is written into \texttt{config.yaml}. Keys absent from the
file fall back to the defaults in \texttt{vae\_model\_kwargs()}
(\texttt{experiments/exp059\_capacity\_freq/codes/exp059\_common.py}, lines 34--55).
\autoref{tab:appb_code_defaults} lists those defaults and marks which of them
\texttt{config.yaml} overrides. \textbf{The twelve rows marked ``code default'' were never
configured for exp059}; they are recorded here because they determine the model that was
built, not because a decision about them is documented anywhere.

\begin{table}[htbp]
\centering
\caption{Defaults in \texttt{vae\_model\_kwargs()} and whether \texttt{config.yaml}
overrides them. ``Code default'' means the key does not appear in any exp059
configuration file.}
\label{tab:appb_code_defaults}
\begin{tabular}{@{}lll@{}}
\hline
Key & Code default & Status in exp059 \\
\hline
\texttt{latent\_dim} & 42 & overridden: 128 \\
\texttt{cond\_dim} & 8 & overridden: 32 \\
\texttt{heatmap\_private\_dim} & 8 & overridden: 40 \\
\texttt{modality\_dropout} & 0.0 & overridden: 0.06 \\
\texttt{freq\_fourier\_features} & 8 & overridden: 16 \\
\texttt{use\_heatmap\_film} & True & overridden: true \\
\texttt{use\_multiscale\_film} & True & code default \\
\texttt{use\_freq\_poe\_expert} & True & code default \\
\texttt{use\_occ\_spatial\_decoder} & True & code default \\
\texttt{use\_occ\_spatial\_tower} & True & code default \\
\texttt{occ\_spatial\_ch} & 8 & code default \\
\texttt{use\_layout\_private\_head} & True & code default \\
\texttt{layout\_private\_hidden} & 384 & code default \\
\texttt{use\_layout\_private\_freq\_film} & True & code default \\
\texttt{graph\_hidden\_dim} & 128 & code default \\
\texttt{graph\_num\_layers} & 3 & code default \\
\texttt{graph\_dropout} & 0.1 & code default \\
\texttt{imp\_peak\_dim} & 8 & code default \\
\hline
\end{tabular}
\end{table}

\texttt{use\_heatmap\_film} is set in \texttt{config.yaml} to the same value as the code
default, so it is listed as overridden but does not change the model.

\section{Fine-tune override configurations}
\label{sec:appb_finetune_configs}

Two override files sit beside the base configuration in
\texttt{experiments/exp059\_capacity\_freq/}: \texttt{config\_al\_finetune.yaml} for the
active-learning fine-tune and \texttt{config\_occ\_only\_finetune.yaml} for the
occupancy-only warm-up. Both are partial: they set a subset of keys and inherit the rest
from \texttt{config.yaml}.

\subsection{Active-learning fine-tune: file value versus value actually used}
\label{subsec:appb_al_finetune}

The fine-tune that runs is not driven directly by \texttt{config\_al\_finetune.yaml}. At
launch, \texttt{write\_runtime\_finetune\_config()}
(\texttt{active\_learning\_pi/al/finetune\_run.py}, lines 61--123) copies that file,
overwrites \texttt{resume\_checkpoint}, recomputes \texttt{num\_epochs} as the checkpoint
epoch plus \texttt{finetune.extra\_epochs}, forces
\texttt{eval\_use\_occ\_only\_layout} to \texttt{True}, and applies the overrides in the
\texttt{finetune} block of the active-learning configuration. The result is written to
\texttt{config\_al\_finetune.runtime.yaml}. \autoref{tab:appb_al_finetune} therefore
prints both columns side by side; where they differ, the runtime column is what ran.

\begin{table}[htbp]
\centering
\small
\caption{Active-learning fine-tune configuration: hand-written
\texttt{config\_al\_finetune.yaml} (29 keys) against the generated
\texttt{config\_al\_finetune.runtime.yaml} (29 keys). Rows in which the two disagree are
discussed in \autoref{sec:appb_conflicts}.}
\label{tab:appb_al_finetune}
\begin{tabular}{@{}lp{0.24\textwidth}p{0.24\textwidth}@{}}
\hline
Key & \texttt{config\_al\_finetune.yaml} & \texttt{...runtime.yaml} \\
\hline
\texttt{resume\_checkpoint} & \texttt{experiments/\allowbreak exp059\_capacity\_freq/\allowbreak checkpoints/\allowbreak last\_model.pt} & \texttt{experiments/\allowbreak exp059\_capacity\_freq/\allowbreak checkpoints/\allowbreak last\_model.pt} \\
\texttt{num\_epochs} & 1050 & 510 \\
\texttt{learning\_rate} & 1.5e-05 & 1.5e-05 \\
\texttt{reset\_lr\_on\_resume} & true & true \\
\texttt{train\_samples\_per\_epoch} & 64000 & 64000 \\
\texttt{al\_overlay\_data\_dir} & \texttt{datasets/\allowbreak data\_multifreq\_\allowbreak al\_overlay\_exp059} & \texttt{datasets/\allowbreak data\_multifreq\_\allowbreak al\_overlay\_exp059} \\
\texttt{al\_overlay\_sample\_weight} & 128.0 & 128.0 \\
\texttt{layout\_train\_prob} & 0.9 & 0.65 \\
\texttt{occ\_only\_encode\_prob} & 0.7 & 0.4 \\
\texttt{eval\_use\_occ\_only\_layout} & true & true \\
\texttt{latent\_distill\_weight} & 2.5 & 2.5 \\
\texttt{output\_distill\_weight} & 2.5 & 2.5 \\
\texttt{use\_beta\_annealing} & false & false \\
\texttt{beta\_final} & 0.03 & 0.03 \\
\texttt{free\_bits} & 0.12 & 0.12 \\
\texttt{heatmap\_weight} & 5.0 & 5.0 \\
\texttt{heatmap\_focus\_heatmap\_weight} & 7.0 & 7.0 \\
\texttt{heatmap\_focus\_start\_frac} & 0.0 & 0.0 \\
\texttt{heatmap\_focus\_start\_epoch} & 0 & 0 \\
\texttt{cross\_freq\_weight} & 1.2 & 1.2 \\
\texttt{eval\_off\_anchor\_mhz} & \texttt{[ 155, 250, 265 ]} & \texttt{[ 155, 250, 265 ]} \\
\texttt{eval\_off\_anchor\_interval} & 5 & 5 \\
\texttt{eval\_off\_anchor\_max\_batches} & 20 & 20 \\
\texttt{al\_finetune\_early\_stop} & true & true \\
\texttt{al\_finetune\_early\_stop\_patience} & 3 & 3 \\
\texttt{al\_finetune\_early\_stop\_min\_delta} & 0.01 & 0.01 \\
\texttt{al\_finetune\_early\_stop\_kind} & \texttt{"layout\_cross"} & \texttt{"layout\_cross"} \\
\texttt{checkpoint\_interval} & 5 & 5 \\
\texttt{epoch\_print\_interval} & 1 & 1 \\
\hline
\end{tabular}
\end{table}

The override generator also copies \texttt{beta\_final}, \texttt{free\_bits} and
\texttt{use\_beta\_annealing} from the occupancy-only warm-up configuration into the
runtime fine-tune configuration (\texttt{finetune\_run.py}, lines 90--92). The keys that
the active-learning configuration's \texttt{finetune} block is allowed to override
(\texttt{finetune\_run.py}, lines 94--117) are: \texttt{al\_overlay\_sample\_weight},
\texttt{layout\_train\_prob}, \texttt{occ\_only\_encode\_prob},
\texttt{eval\_use\_occ\_only\_layout}, \texttt{latent\_distill\_weight},
\texttt{output\_distill\_weight}, \texttt{latent\_distill\_logvar\_weight},
\texttt{heatmap\_phys\_p99\_weight}, \texttt{eval\_off\_anchor\_mhz},
\texttt{eval\_off\_anchor\_mhz\_weights}, \texttt{eval\_off\_anchor\_interval},
\texttt{eval\_off\_anchor\_max\_batches}, \texttt{al\_finetune\_early\_stop},
\texttt{al\_finetune\_early\_stop\_patience},
\texttt{al\_finetune\_early\_stop\_min\_delta},
\texttt{al\_finetune\_early\_stop\_kind}, \texttt{checkpoint\_interval},
\texttt{beta\_final}, \texttt{free\_bits} and \texttt{use\_beta\_annealing}.

\subsection{Occupancy-only warm-up configuration}
\label{subsec:appb_occ_warmup}

\begin{table}[htbp]
\centering
\small
\caption{\texttt{config\_occ\_only\_finetune.yaml}, the occupancy-only warm-up override
(28 keys).}
\label{tab:appb_occ_warmup}
\begin{tabular}{@{}lp{0.36\textwidth}@{}}
\hline
Key & Value \\
\hline
\texttt{resume\_checkpoint} & \texttt{experiments/\allowbreak exp059\_capacity\_freq/\allowbreak checkpoints/\allowbreak last\_model.pt} \\
\texttt{num\_epochs} & 1100 \\
\texttt{learning\_rate} & 1.2e-05 \\
\texttt{reset\_lr\_on\_resume} & true \\
\texttt{train\_samples\_per\_epoch} & 56000 \\
\texttt{layout\_train\_prob} & 1.0 \\
\texttt{occ\_only\_encode\_prob} & 1.0 \\
\texttt{eval\_use\_occ\_only\_layout} & true \\
\texttt{use\_beta\_annealing} & false \\
\texttt{beta\_final} & 0.03 \\
\texttt{free\_bits} & 0.12 \\
\texttt{latent\_distill\_weight} & 2.5 \\
\texttt{output\_distill\_weight} & 2.5 \\
\texttt{cross\_freq\_weight} & 1.2 \\
\texttt{cross\_freq\_layout\_z\_only} & true \\
\texttt{heatmap\_weight} & 5.0 \\
\texttt{heatmap\_focus\_heatmap\_weight} & 6.0 \\
\texttt{heatmap\_focus\_start\_frac} & 0.0 \\
\texttt{heatmap\_focus\_start\_epoch} & 0 \\
\texttt{eval\_off\_anchor\_mhz} & \texttt{[ 155, 250, 265 ]} \\
\texttt{eval\_off\_anchor\_interval} & 5 \\
\texttt{eval\_off\_anchor\_max\_batches} & 20 \\
\texttt{al\_finetune\_early\_stop} & true \\
\texttt{al\_finetune\_early\_stop\_patience} & 4 \\
\texttt{al\_finetune\_early\_stop\_min\_delta} & 0.01 \\
\texttt{al\_finetune\_early\_stop\_kind} & \texttt{"layout\_cross"} \\
\texttt{checkpoint\_interval} & 5 \\
\texttt{epoch\_print\_interval} & 1 \\
\hline
\end{tabular}
\end{table}

The warm-up sets both \texttt{layout\_train\_prob} and \texttt{occ\_only\_encode\_prob}
to $1.0$, i.e.\ it trains the occupancy-only encode path exclusively.

\subsection{Active-learning configuration blocks that reach the fine-tune}
\label{subsec:appb_al_blocks}

\begin{table}[htbp]
\centering
\small
\caption{\texttt{finetune} block of the primary active-learning configuration
\texttt{active\_learning\_pi/config/exp059\_gp\_error.json} (lines 111--122).}
\label{tab:appb_al_finetune_block}
\begin{tabular}{@{}lp{0.38\textwidth}@{}}
\hline
Key & Value \\
\hline
\texttt{enabled} & true \\
\texttt{run\_after\_cycle} & false \\
\texttt{experiment\_dir} & \texttt{experiments/\allowbreak exp059\_capacity\_freq} \\
\texttt{config\_path} & \texttt{experiments/\allowbreak exp059\_capacity\_freq/\allowbreak config\_al\_finetune.yaml} \\
\texttt{checkpoint\_path} & \texttt{experiments/\allowbreak exp059\_capacity\_freq/\allowbreak checkpoints/\allowbreak last\_model.pt} \\
\texttt{extra\_epochs} & 45 \\
\texttt{overlay\_data\_dir} & \texttt{datasets/\allowbreak data\_multifreq\_\allowbreak al\_overlay\_exp059} \\
\texttt{al\_overlay\_sample\_weight} & 128.0 \\
\texttt{occ\_only\_encode\_prob} & 0.4 \\
\texttt{layout\_train\_prob} & 0.65 \\
\hline
\end{tabular}
\end{table}

\begin{table}[htbp]
\centering
\caption{\texttt{pre\_al\_occ\_only\_warmup} block of
\texttt{active\_learning\_pi/config/exp059\_gp\_error.json} (lines 106--110).}
\label{tab:appb_al_warmup_block}
\begin{tabular}{@{}lp{0.38\textwidth}@{}}
\hline
Key & Value \\
\hline
\texttt{config\_path} & \texttt{experiments/\allowbreak exp059\_capacity\_freq/\allowbreak config\_occ\_only\_finetune.yaml} \\
\texttt{checkpoint\_path} & \texttt{experiments/\allowbreak exp059\_capacity\_freq/\allowbreak checkpoints/\allowbreak last\_model.pt} \\
\texttt{extra\_epochs} & 100 \\
\hline
\end{tabular}
\end{table}

The acquisition surrogate has its own settings in the \texttt{gp\_error} block of the same
file, transcribed in \autoref{tab:appb_gp_error}. The primary active-learning
configuration sets \texttt{acquisition\_mode} to \texttt{"gp\_error"} (line 139); within
the block the GP type is \texttt{"svgp"}, the fitted target is \texttt{"p99\_ae"} on the
\texttt{"pred"} path, and \texttt{score\_mode} is \texttt{"mu"}.
\fillin{state what \texttt{kappa}, \texttt{novelty\_weight}, \texttt{novelty\_k},
\texttt{joint\_alpha\_hm} and \texttt{joint\_alpha\_imp} do to the acquisition score, and
in particular which of them are inactive when \texttt{score\_mode} is \texttt{"mu"} and
\texttt{novelty\_weight} is $0.0$; the scoring function in
\texttt{active\_learning\_pi/al/gp\_error\_surrogate.py} would settle it.}

\begin{table}[htbp]
\centering
\caption{\texttt{gp\_error} block of
\texttt{active\_learning\_pi/config/exp059\_gp\_error.json} (lines 140--155), 14 keys.}
\label{tab:appb_gp_error}
\begin{tabular}{@{}lp{0.30\textwidth}@{}}
\hline
Key & Value \\
\hline
\texttt{path} & \texttt{"pred"} \\
\texttt{target} & \texttt{"p99\_ae"} \\
\texttt{gp} & \texttt{"svgp"} \\
\texttt{kappa} & 0.5 \\
\texttt{n\_fit} & 4000 \\
\texttt{n\_inducing} & 256 \\
\texttt{svgp\_epochs} & 80 \\
\texttt{seed} & 0 \\
\texttt{verbose} & true \\
\texttt{novelty\_weight} & 0.0 \\
\texttt{novelty\_k} & 10 \\
\texttt{joint\_alpha\_hm} & 1.0 \\
\texttt{joint\_alpha\_imp} & 1.0 \\
\texttt{score\_mode} & \texttt{"mu"} \\
\hline
\end{tabular}
\end{table}

Three settings that a sparse variational GP normally needs are genuinely absent from this
block: the minibatch size used to fit the variational parameters, the learning rate, and
the kernel.
\fillin{state the residual GP's batch size, learning rate and kernel --- they are not set
in \texttt{active\_learning\_pi/config/exp059\_gp\_error.json}; if they are library
defaults or hard-coded in \texttt{active\_learning\_pi/al/gp\_error\_surrogate.py}, name
the source, and for a library default name the library and version that supplies it.}

\clearpage

\section{Conflicts between configuration files}
\label{sec:appb_conflicts}

Four disagreements between configuration artifacts were found while assembling this
appendix. They are recorded here rather than resolved, because in each case two files that
are both in use state different things and no authoritative artifact settles which was
intended.

\subsection*{C1: fine-tune layout probabilities}

\texttt{experiments/exp059\_capacity\_freq/config\_al\_finetune.yaml} (lines 11--12) sets
\texttt{layout\_train\_prob} $=0.9$ and \texttt{occ\_only\_encode\_prob} $=0.7$. The same
pair appears in \texttt{docs/cursor-handoff.md} (lines 42 and 77),
\texttt{docs/gp-error-surrogate.md} (line 86) and \texttt{CLAUDE.md} (line 17).

However, \texttt{active\_learning\_pi/config/exp059\_gp\_error.json} (lines 120--121)
overrides the pair in its \texttt{finetune} block to \texttt{layout\_train\_prob} $=0.65$
and \texttt{occ\_only\_encode\_prob} $=0.4$. \texttt{write\_runtime\_finetune\_config()}
(\texttt{finetune\_run.py}, lines 94--117) lets the active-learning JSON win over the
hand-written fine-tune file for both keys, and the generated
\texttt{config\_al\_finetune.runtime.yaml} (lines 9--10) accordingly carries $0.65$ and
$0.4$. A fine-tune launched from the current configuration therefore runs at $0.65/0.4$,
not at the documented $0.9/0.7$. Both values are stated in
\autoref{tab:appb_al_finetune} and \autoref{tab:appb_al_finetune_block}; neither is
presented here as the correct one.

What the four completed active-learning cycles used is a separate question, and it cannot
be answered from the repository. The runtime file is rewritten at every fine-tune launch,
so it records only the most recent one, and no per-cycle configuration record survives: no
iteration directory under
\texttt{active\_learning\_pi/runs/al\_exp059\_gp\_error\_001/} records
\texttt{layout\_train\_prob} at all, and neither \texttt{state.json} nor
\texttt{decision\_ledger.json} carries these keys. This is a reproducibility gap, not a
detail: the two probabilities govern how much of the fine-tune trains the occupancy-only
encode path, which is the path the active-learning loop scores with.
\fillin{pin down, from surviving training logs, which \texttt{layout\_train\_prob} and
\texttt{occ\_only\_encode\_prob} were in force in each of the four completed cycles of
\texttt{al\_exp059\_gp\_error\_001}. The repository preserves no per-cycle configuration
snapshot, so a fine-tune console log or a checkpoint-side metadata record is the only
thing that can settle it.}
\fillin{decide which pair was intended, correct whichever file is wrong, and --- once the
per-cycle values are recovered --- state in Chapter 5 whether any reported active-learning
result is affected by the value actually used.}
\fillin{if the per-cycle values cannot be recovered, add a per-iteration configuration
snapshot to the active-learning pipeline so that future cycles are reproducible, and say
so in the limitations section.}

\subsection*{C2: fine-tune epoch count}

\texttt{config\_al\_finetune.yaml} (line 5) states \texttt{num\_epochs} $=1050$, while the
generated \texttt{config\_al\_finetune.runtime.yaml} (line 3) states $510$.
\texttt{finetune\_run.py} (lines 71--79) recomputes \texttt{num\_epochs} as the checkpoint
epoch plus \texttt{extra\_epochs}, which is $45$ in \texttt{exp059\_gp\_error.json}
(line 117). The value in the hand-written file is therefore never the number of epochs
that runs. The base \texttt{config.yaml} independently holds \texttt{num\_epochs} $=510$.
\fillin{state the checkpoint epoch that the runtime recomputation started from, from the
checkpoint metadata or the fine-tune log, so the 510 in the runtime file can be traced.}

\subsection*{C3: stale experiment name in the override headers}

The header comments on line 1 of both
\texttt{experiments/exp059\_capacity\_freq/config\_al\_finetune.yaml} and
\texttt{experiments/exp059\_capacity\_freq/config\_occ\_only\_finetune.yaml} read
``for exp058'', while both files live under \texttt{exp059\_capacity\_freq} and set
\texttt{resume\_checkpoint} (line 4 in each) to
\texttt{experiments/exp059\_capacity\_freq/checkpoints/last\_model.pt}. The
\texttt{config\_al\_finetune.yaml} header additionally reads ``only after fresh train has
produced last\_model.pt / Do NOT point at exp057 checkpoints.'' The comment lines are
stripped by \texttt{load\_yaml\_config()} before parsing, so they do not affect execution,
but they mislabel the file.

\subsection*{C4: normalisation statistics path}

\texttt{experiments/exp059\_capacity\_freq/config.yaml} (line 134) points the training set
at \texttt{datasets/data\_multifreq\_train\_norm\_unbounded}, whereas
\texttt{pipelines/latent/optimize.py} (line 105) sets
\texttt{NORMALIZATION\_STATS\_PATH} to
\texttt{datasets/data\_multifreq\_norm/normalization\_stats.json}. The two are different
directories.
\fillin{state whether the two normalisation-statistics files contain the same per-MHz
statistics; if they do not, any comparison between a Stage-2 result and an exp059
prediction is in different units and must be flagged in Chapter 4.}

\subsection*{C5: the base configuration is not the capacity-phase record}

Beyond the four file-to-file conflicts, \texttt{config.yaml} has been edited since the
capacity-phase run. \texttt{docs/cursor-handoff.md} states that the capacity phase ran
with \texttt{layout\_train\_prob} $=0$ and \texttt{occ\_only\_encode\_prob} $=0$, and the
file now holds $0.65$ and $0.4$ (\autoref{tab:appb_crossfreq}). The current file therefore
does not describe the run that produced the capacity checkpoint. This is stated plainly
because the tables in \autoref{sec:appb_exp059_config} would otherwise be read as a run
record, which they are not.
\fillin{check the remaining 181 keys against the training log or a config snapshot taken
at the capacity run, and mark in \autoref{sec:appb_exp059_config} every key whose current
value differs from the value used then.}

\section{Stage-2 latent-optimisation configuration}
\label{sec:appb_stage2}

\textbf{The configuration in this section is not exp059's.}
\texttt{pipelines/latent/optimize.py} still selects \texttt{EXPERIMENT =
"exp038\_true\_multi"} for the occupancy decoder and the impedance surrogate (lines
45--54), so no Stage-2 hyperparameter set wired to \texttt{exp059\_capacity\_freq} exists.
The tables below record the \texttt{CONFIG} block as it stands: an exp038-targeted
configuration. They must not be presented as the settings under which an exp059 model was
inverted.

Like every runnable script in this repository, \texttt{optimize.py} takes no command-line
arguments; the values are module-level constants edited before the run. Several are
overridable through environment variables, which is noted per row.

\begin{table}[htbp]
\centering
\caption{\texttt{optimize.py} model and checkpoint selection (lines 45--54).}
\label{tab:appb_opt_model}
\begin{tabular}{@{}lp{0.42\textwidth}@{}}
\hline
Constant & Value \\
\hline
\texttt{EXPERIMENT} & \texttt{"exp038\_true\_multi"} \\
\texttt{CHECKPOINT\_PATH} & \texttt{f"experiments/\{EXPERIMENT\}/\allowbreak checkpoints/last\_model.pt"} \\
\texttt{SURROGATE\_CHECKPOINT\_PATH} & \texttt{f"experiments/\{EXPERIMENT\}/\allowbreak checkpoints/surrogate\_best.pt"} \\
\texttt{USE\_SURROGATE} & \texttt{\_env\_flag("LATENT\_OPT\_USE\_SURROGATE", "1")} \\
\hline
\end{tabular}
\end{table}

\begin{table}[htbp]
\centering
\caption{\texttt{optimize.py} search setup (lines 56--69).}
\label{tab:appb_opt_search}
\begin{tabular}{@{}lp{0.42\textwidth}@{}}
\hline
Constant & Value \\
\hline
\texttt{K\_LIST} & \texttt{list(range(1, 26))} \\
\texttt{NUM\_CANDIDATE\_SEEDS} & 32 \\
\texttt{SEED\_MODE} & \texttt{os.getenv("LATENT\_OPT\_SEED\_MODE", "random").strip().lower()} \\
\texttt{MASTER\_SEED} & \texttt{\_parse\_master\_seed()} (\texttt{None} when \texttt{LATENT\_OPT\_MASTER\_SEED} is unset) \\
\texttt{OPTIMIZE\_BATCH\_PER\_K} & True \\
\hline
\end{tabular}
\end{table}

\begin{table}[htbp]
\centering
\caption{\texttt{optimize.py} latent-optimiser settings and objective (lines 71--77).}
\label{tab:appb_opt_optimizer}
\begin{tabular}{@{}lp{0.40\textwidth}@{}}
\hline
Constant & Value \\
\hline
\texttt{NUM\_STEPS} & 1200 \\
\texttt{LR} & 5e-2 \\
\texttt{GRAD\_CLIP} & 5.0 \\
\texttt{INIT\_SHARED\_TEMP} & 0.8 \\
\texttt{OBJECTIVE\_MODE} & \texttt{"gap\_max"} (alternative: \texttt{"feasible\_map"}) \\
\texttt{LOSS\_SPACE} & \texttt{"log"} (alternative: \texttt{"ohm"}) \\
\hline
\end{tabular}
\end{table}

The optimiser over the latent vector is \texttt{torch.optim.Adam([z], lr=LR)}
(\texttt{optimize.py}, lines 504 and 514), with everything else in the model frozen. The
discrete read-out from continuous occupancy probabilities to a hard top-$K$ placement is
\texttt{\_ste\_topk(occ\_prob, K)} (\texttt{optimize.py}, line 259), the straight-through
estimator that keeps the gradient path to $\mathbf{z}$ intact.

\begin{table}[htbp]
\centering
\caption{\texttt{optimize.py} latent-prior and spectrum-shape regularisation weights
(lines 79--86).}
\label{tab:appb_opt_prior}
\begin{tabular}{@{}ll@{}}
\hline
Constant & Value \\
\hline
\texttt{Z\_L2\_WEIGHT} & 0.05 \\
\texttt{Z\_PRIOR\_MODE} & \texttt{"agg\_posterior"} \\
\texttt{SHAPE\_REG\_WEIGHT} & 2.0 \\
\texttt{SHAPE\_MIN\_STD\_DLOG} & 0.12 \\
\texttt{SHAPE\_TRACK\_WEIGHT} & 2.0 \\
\texttt{POSTERIOR\_BOUNDARY\_WEIGHT} & 10.0 \\
\texttt{POSTERIOR\_SIGMA\_LIMIT} & 2.5 \\
\texttt{PHYSICS\_AR\_WEIGHT} & 0.5 \\
\hline
\end{tabular}
\end{table}

\begin{table}[htbp]
\centering
\caption{\texttt{optimize.py} peak-aware spectrum terms (lines 88--94). Each is read from
the named environment variable with the default shown.}
\label{tab:appb_opt_peak}
\begin{tabular}{@{}lll@{}}
\hline
Constant & Environment variable & Default \\
\hline
\texttt{PEAK\_LOSS\_WEIGHT} & \texttt{LATENT\_OPT\_PEAK\_LOSS\_WEIGHT} & \texttt{"1.0"} \\
\texttt{PEAK\_INDEX\_WEIGHT} & \texttt{LATENT\_OPT\_PEAK\_INDEX\_WEIGHT} & \texttt{"3.0"} \\
\texttt{PEAK\_MAG\_WEIGHT} & \texttt{LATENT\_OPT\_PEAK\_MAG\_WEIGHT} & \texttt{"2.0"} \\
\texttt{DUAL\_TOPK\_WEIGHT} & \texttt{LATENT\_OPT\_DUAL\_TOPK\_WEIGHT} & \texttt{"2.5"} \\
\texttt{IMPEDANCE\_NUM\_PEAKS} & \texttt{LATENT\_OPT\_IMPEDANCE\_NUM\_PEAKS} & \texttt{"8"} \\
\texttt{IMPEDANCE\_TOPK\_K} & \texttt{LATENT\_OPT\_IMPEDANCE\_TOPK\_K} & \texttt{"20"} \\
\hline
\end{tabular}
\end{table}

The source comment at line 88 records that these terms are ``aligned with exp038
impedance\_spectrum\_loss'', which is consistent with the exp038 target of the whole
block.

\begin{table}[htbp]
\centering
\caption{\texttt{optimize.py} diversity, occupancy-confidence and feasibility terms
(lines 96--103).}
\label{tab:appb_opt_feasibility}
\begin{tabular}{@{}ll@{}}
\hline
Constant & Value \\
\hline
\texttt{DIVERSITY\_WEIGHT} & 0.35 \\
\texttt{OCC\_CONFIDENCE\_WEIGHT} & 0.5 \\
\texttt{BOUNDARY\_MARGIN} & 0.1 \\
\texttt{GAP\_REWARD\_WEIGHT} & 1.0 \\
\texttt{EXCEED\_WEIGHT} & 25.0 \\
\texttt{EXCEED\_POWER} & 2.0 \\
\texttt{SELECT\_METRIC} & \texttt{"max\_ohm"} \\
\hline
\end{tabular}
\end{table}

\begin{table}[htbp]
\centering
\caption{\texttt{optimize.py} paths, device and output naming (lines 105--120).}
\label{tab:appb_opt_paths}
\begin{tabular}{@{}lp{0.44\textwidth}@{}}
\hline
Constant & Value \\
\hline
\texttt{NORMALIZATION\_STATS\_PATH} & \texttt{datasets/\allowbreak data\_multifreq\_norm/\allowbreak normalization\_stats.json} \\
\texttt{TARGET\_IMPEDANCE\_PATH} & \texttt{configs/\allowbreak target\_impedance.npy} \\
\texttt{OUTPUT\_ROOT} & \texttt{data/latent\_runs} \\
\texttt{DEVICE} & \texttt{"cuda" if torch.cuda.is\_available() else "cpu"} \\
\texttt{DTYPE} & \texttt{torch.float32} \\
\texttt{HIST\_KEYS} & \texttt{("total", "exceed", "gap\_mean", "best\_score")} \\
\texttt{SOLUTION\_PREFIX} & \texttt{"best\_"} \\
\texttt{SOLUTION\_LATENT\_NAME} & \texttt{f"\{SOLUTION\_PREFIX\}latent.npy"} \\
\texttt{SOLUTION\_METRICS\_NAME} & \texttt{f"\{SOLUTION\_PREFIX\}metrics.json"} \\
\hline
\end{tabular}
\end{table}

\texttt{NORMALIZATION\_STATS\_PATH} is the second half of conflict C4 in
\autoref{sec:appb_conflicts}: it does not point at the exp059 training set named in
\texttt{config.yaml}.

\clearpage

\section{Values this appendix cannot supply}
\label{sec:appb_missing}

For completeness, the following are known gaps rather than omissions.

\begin{itemize}
  \item No per-key rationale exists for any of the 183 keys in \texttt{config.yaml}; the
        file contains no comment lines.
  \item No sweep record or hyperparameter table was located that distinguishes exp059
        values that were tuned from values inherited from exp057 or exp058.
  \item Twelve architecture switches (\autoref{tab:appb_code_defaults}) have no documented
        value in any exp059 configuration file and exist only as code defaults.
  \item The effective per-GPU and global batch size under DDP, and the learning rate
        produced by \texttt{ddp\_linear\_lr\_scale}, are not recorded in any artifact.
  \item No Stage-2 hyperparameter set targeting exp059 exists;
        \autoref{sec:appb_stage2} documents an exp038-targeted configuration.
  \item The residual-GP acquisition surrogate's minibatch size, learning rate and kernel
        are not set in the \texttt{gp\_error} block (\autoref{tab:appb_gp_error}); the
        remaining fourteen keys of that block are transcribed there.
  \item The \texttt{layout\_train\_prob} and \texttt{occ\_only\_encode\_prob} values in
        force during each completed active-learning cycle are not recoverable from the
        repository: the runtime fine-tune configuration is overwritten at every launch and
        no per-iteration snapshot is written (C1, \autoref{sec:appb_conflicts}).
\end{itemize}

% ---------------------------------------------------------------
% FILL IN checklist:
%  1. sec:appb_scope - provenance of tuned vs inherited exp059 values; check experiments/exp059_capacity_freq/notes.md, no sweep artifact located.
%  2. subsec:appb_split - does train_core.py read k_balance_power / freq_balance_power when balance_k / balance_freq are false? State the sampler actually used.
%  3. subsec:appb_heatmap_loss - move per-term explanation of the heatmap losses into Chapter 3; config records no rationale.
%  4. tab:appb_heatmap_clip - origin of -1.6791870594024658 (shared by heatmap_z_clip_min and physics_fg_clip_min); check normalization_stats.json. Remove the word "presumably" once settled.
%  5. subsec:appb_impedance_loss - which curriculum form train_core.py uses (*_frac vs *_epoch) and the role of recalculate_curriculum_on_resume.
%  6. subsec:appb_kl - which single beta value applies when use_beta_annealing is false.
%  7. subsec:appb_crossfreq - recover the layout_train_prob / occ_only_encode_prob values in force during the capacity-phase run (training log or config snapshot next to checkpoints/last_model.pt).
%  8. subsec:appb_runtime - effective per-GPU/global batch size and the DDP-scaled learning rate, from a training log line.
%  9. subsec:appb_eval_keys - why 250.0 MHz (a training anchor) sits in eval_off_anchor_mhz, and whether off-anchor metrics are reported per frequency or pooled.
% 10. subsec:appb_al_blocks - role of kappa / novelty_weight / novelty_k / joint_alpha_hm / joint_alpha_imp in the acquisition score when score_mode="mu"; check gp_error_surrogate.py.
% 11. subsec:appb_al_blocks - residual-GP batch size, learning rate and kernel: not in exp059_gp_error.json; name gp_error_surrogate.py or the library+version supplying the defaults.
% 12. sec:appb_conflicts C1 - pin the per-cycle layout_train_prob / occ_only_encode_prob for the four completed cycles from surviving training logs (no per-iteration snapshot exists in the repo).
% 13. sec:appb_conflicts C1 - decide intended layout probability pair; then state in Chapter 5 whether any AL result is affected.
% 14. sec:appb_conflicts C1 - add a per-iteration config snapshot to the AL pipeline, or record the gap in the limitations section.
% 15. sec:appb_conflicts C2 - checkpoint epoch that the runtime num_epochs recomputation started from.
% 16. sec:appb_conflicts C4 - whether the two normalisation-statistics files hold the same per-MHz statistics.
% 17. sec:appb_conflicts C5 - diff the remaining 181 keys against the capacity-run snapshot and mark every key that changed.
%
% Citation verification queue:
%  - None. This appendix cites no literature by design; every statement is sourced to a
%    configuration file or a source-code line range given in the evidence block.
%
% Evidence warnings:
%  - C1 (layout_train_prob / occ_only_encode_prob 0.9/0.7 vs 0.65/0.4): documented value
%    and the value the current configuration would execute differ. Both printed in
%    tab:appb_al_finetune; not resolved here. What the four COMPLETED cycles of
%    al_exp059_gp_error_001 actually ran at is NOT recoverable from the repository - the
%    runtime config is overwritten at every launch and no iteration directory, state.json
%    or decision_ledger.json records these keys. Stated in the text as a reproducibility
%    gap. Do not write "the completed cycles ran at 0.65/0.4" anywhere in the thesis.
%  - C2 (num_epochs 1050 in config_al_finetune.yaml vs 510 in the generated runtime file;
%    config.yaml also 510). Not resolved. The runtime value is recomputed at launch, so
%    the base file value never runs. No arithmetic reconstruction of the checkpoint epoch
%    was performed here - deliberately, since that would be recomputation.
%  - C3 (both override files carry "for exp058" headers while pointing at exp059
%    checkpoints). Cosmetic at runtime because comment lines are stripped, but mislabels
%    the artifact.
%  - C4 (normalisation stats path differs between config.yaml and optimize.py).
%  - C5 (config.yaml no longer describes the capacity-phase run per docs/cursor-handoff.md:
%    capacity phase used layout_train_prob=0 / occ_only_encode_prob=0). Tables in
%    sec:appb_exp059_config must not be read as a run record.
%  - Stage 2 (sec:appb_stage2) targets exp038_true_multi. No exp059 optimiser
%    hyperparameter set exists; nothing in this appendix may be used to attribute an
%    end-to-end inverse-design result to exp059.
%  - Table grouping into 21 groups is editorial, not structure present in config.yaml; the
%    cited line ranges are spans, and the keys inside a group are not contiguous in the file.
%  - eval_off_anchor_mhz = [155.0, 250.0, 265.0] is NOT a held-out set: 250.0 is one of
%    the 24 training anchors (datasets/data_multifreq_train_norm_unbounded/dataset_meta.json,
%    pi_frequencies_mhz) and carries the highest weight (3.0 vs 2.5, 2.5). The metric is
%    also a checkpoint-selection criterion - checkpoints/best_off_anchor_model.pt exists.
%    Any frequency-generalisation claim must be restricted to 155.0 and 265.0 MHz.
%  - The gp_error block (tab:appb_gp_error, 14 keys, lines 140-155) was read directly from
%    active_learning_pi/config/exp059_gp_error.json during this revision, not taken from
%    the evidence block. Only batch size, learning rate and kernel are genuinely absent.
%  - No experimental result, ranking, or improvement claim appears anywhere in this
%    appendix, by design.
% ---------------------------------------------------------------
